% SOURCE_STATEMENT: "Problem 22. Assume E = 0 and angular momentum = 0. Can Jacobi's metric give new insights beyond what McGehee's coordinates give? (Richard Montgomery)"
% Preparation and final self-audit: CONTRIBUTING.md prompts C and D; source/audit artifacts retained in the local physics_statements_2026_10_11 batch.
% !TeX program = lualatex
\documentclass[11pt]{article}
\usepackage[margin=25mm]{geometry}
\usepackage{fontspec}
\setmainfont{Latin Modern Roman}
\usepackage{amsmath,amssymb,mathtools}
\usepackage{unicode-math}
\setmathfont{Latin Modern Math}
\usepackage{xurl,hyperref}
\hypersetup{colorlinks=true,urlcolor=blue}
\setcounter{secnumdepth}{0}
\setlength{\parindent}{0pt}
\setlength{\parskip}{5pt}
\setlength{\emergencystretch}{3em}
\newcommand{\sref}[2]{\hyperlink{source:#1:#2}{[S#2]}}
\newcommand{\cataloguescope}[1]{\paragraph{Recorded target.}#1}
\newcommand{\R}{\mathbb R}
\newcommand{\C}{\mathbb C}
\newcommand{\N}{\mathbb N}
\newcommand{\Z}{\mathbb Z}
\newcommand{\Q}{\mathbb Q}
\newcommand{\F}{\mathbb F}
\newcommand{\eps}{\varepsilon}
\begin{document}
% UnsolvedMath ID 20002551; source catalogue code AIM-PHYSICS-0059
\setcounter{section}{20002550}
\section{A Jacobi-Maupertuis and McGehee scale-shape dictionary}\label{20002551}
{\small UnsolvedMath ID 20002551 · Mathematical Physics}
\subsection{Definitions and mathematical statement}

\paragraph{Shared notation.}$\mathbb N=\{1,2,\ldots\}$; the standard number fields have their usual meanings. All additional conventions are stated locally.


\paragraph{Definitions and notation.}In the Newtonian $n$-body problem, positions $q_i\in\mathbb R^d$ and masses $m_i>0$ obey $m_i\ddot q_i=\nabla_{q_i}U$, where $U(q)=\sum_{i<j}m_im_j/|q_i-q_j|$ and the gravitational constant is normalized to one. Here $T=\tfrac12\sum_i m_i|\dot q_i|^2$, $E=T-U$, and a collision has $q_i=q_j$. The center of mass can be fixed at zero; moment of inertia is $I=\sum_i m_i|q_i|^2$. At fixed energy $E$, the Jacobi--Maupertuis metric on the Hill region is conformal to the mass metric by the factor $2(E+U)$. Its geodesics, with appropriate reparametrization away from degeneracies, describe mechanical trajectories. McGehee coordinates rescale position and time to analyze collision or asymptotic boundaries. The two frameworks encode different geometric features.

\paragraph{Catalogue wording retained for traceability.}
Assume $E=0$ and angular momentum zero in the source’s celestial-mechanics setting. Can Jacobi’s metric give insights beyond those obtained from McGehee coordinates?
\paragraph{Scope and source qualification.}This is a request for useful mathematical insight, not a fixed theorem or an asserted equivalence of coordinate constructions. The conditions E=0 and angular momentum=0 remain explicit. The catalogue title/status was paired with a machine-generated research attempt; that attempt is not peer-reviewed evidence of settlement. The statement covers the original source questions rather than only the title’s subcase.
\subsection{Short English statement}
Determine whether the Jacobi metric yields new zero-energy, zero-angular-momentum insights beyond McGehee coordinates.
\subsection{Sources}
\begin{itemize}
\item[\hypertarget{source:20002551:1}{S1}] ULAM.AI and the UnsolvedMath Contributors, UnsolvedMath, record 20002551 (AIM-PHYSICS-0059). Catalogue curation under CC BY 4.0; definitions and English exposition added locally.
\url{https://huggingface.co/datasets/ulamai/UnsolvedMath}.
\item[\hypertarget{source:20002551:2}{S2}] American Institute of Mathematics, Variational Methods in Celestial Mechanics; notes by Kuo-Chang Chen, version 22 June 2003, source item 22.
\url{https://aimath.org/WWN/varcelest/varcelest.pdf}.
\end{itemize}
\label{end:20002551}
\end{document}
